\chapter{Bloco 3 -- Relações lógicas}
\label{cap:bloco3}

Associamos a cada proposição simples $P_i$ uma binária $\delta_i$ ($\delta_i=1\iff P_i$ verdadeira). Proposições
compostas viram \textbf{restrições lineares sobre os $\delta_i$}. Releia a Seção~\ref{sec:logica} antes de
começar: lá está a técnica de ``proibir as linhas falsas'' da tabela-verdade.

\section{3A -- Laboratório da tabela-verdade}

\paragraph{Como o verificador funciona.} Cada desafio é uma tupla
\py{(nome, nº de variáveis, regra lógica, restrição linear)}. A função \py{verificar_logicas} testa \textbf{todas} as
$2^n$ atribuições de 0/1 e compara: a restrição deve ser satisfeita \emph{exatamente} nas linhas em que a regra é
verdadeira. Você escreve a restrição como uma \py{lambda} que devolve \py{True}/\py{False}; para combinar várias
restrições, use \py{and}.

\begin{discussao}
\begin{itemize}
  \item \emph{$\delta_1+\delta_2\ge1$ representa qual conectivo?} O ``ou inclusivo'': só a linha $(0,0)$ viola.
  \emph{E $\delta_1-\delta_2\le0$?} A implicação $P_1\Rightarrow P_2$: só a linha $(1,0)$ viola.
  \item \emph{Como representar $P_1\Rightarrow P_2$?} Ela só é falsa em $P_1=1, P_2=0$. Precisamos de uma restrição
  violada só nessa linha: $\delta_1\le\delta_2$.
\end{itemize}
\end{discussao}

\subsection*{Derivando cada linha do exercício}

\begin{center}
\small
\begin{tabularx}{\linewidth}{lXl}
\toprule
Regra & Raciocínio (linhas falsas $\to$ restrição) & \py{lambda}\\
\midrule
$\neg P_1$ & falsa em $\delta_1=1$ & \py{a == 0}\\
$P_1\wedge P_2$ & verdadeira só em $(1,1)$: a soma tem que ser 2 & \py{a + b == 2}\\
$P_1\vee P_2$ & falsa só em $(0,0)$: soma $\ge1$ & \py{a + b >= 1}\\
$P_1\xor P_2$ & falsa em $(0,0)$ e $(1,1)$: soma exatamente 1 & \py{a + b == 1}\\
$P_1\Rightarrow P_2$ & $\neg P_1\vee P_2$: $(1-a)+b\ge1$ & \py{a - b <= 0}\\
$P_1\Leftrightarrow P_2$ & falsa em $(1,0)$ e $(0,1)$: iguais & \py{a - b == 0}\\
$P_1\wedge P_2\Rightarrow P_3$ & falsa só em $(1,1,0)$; fórmula geral com $k=2$ & \py{a + b - c <= 1}\\
$P_1\Rightarrow P_2\vee P_3$ & falsa só em $(1,0,0)$; fórmula geral com $k=1$ & \py{a - b - c <= 0}\\
pelo menos 2 de 4 & contar os verdadeiros & \py{a + b + c + d >= 2}\\
exatamente 2 de 4 & idem & \py{a + b + c + d == 2}\\
no máximo 1 de 4 & exclusão mútua & \py{a + b + c + d <= 1}\\
\bottomrule
\end{tabularx}
\end{center}

\begin{exemplo}[conferindo $P_1\wedge P_2\Rightarrow P_3$ pela tabela]
\begin{center}
\begin{tabular}{ccc|c|c|c}
\toprule
$\delta_1$ & $\delta_2$ & $\delta_3$ & regra & $\delta_1+\delta_2-\delta_3$ & $\le1$?\\
\midrule
1&1&1& V & 1 & sim\\
1&1&0& \textbf{F} & \textbf{2} & \textbf{não}\\
1&0&$*$& V & $\le1$ & sim\\
0&1&$*$& V & $\le1$ & sim\\
0&0&$*$& V & $\le0$ & sim\\
\bottomrule
\end{tabular}
\end{center}
A restrição é violada exatamente na única linha falsa \checkmark.
\end{exemplo}

\begin{atencao}[o verificador aceita, mas o solver não]
O verificador avalia expressões Python sobre números 0/1, então aceita coisas como \py{a + b > 0} ou
\py{a - b != 0}, que dão o resultado certo na tabela. Mas \textbf{nenhum solver de PL/PI aceita \py{>} estrito nem
\py{!=}}. Treine escrever a versão que funcionaria no PuLP: \py{a + b >= 1} e \py{a + b == 1}.
\end{atencao}

\section{3B -- Carteira de investimentos}

\paragraph{O problema.} Cinco investimentos ($x_j=1$ se investir em $j$), orçamento 14.
\begin{center}
\begin{tabular}{lccccc}
\toprule
investimento $j$ & 1 & 2 & 3 & 4 & 5\\
\midrule
retorno & 5 & 6 & 4 & 10 & 10\\
custo & 4 & 6 & 5 & 3 & 5\\
\bottomrule
\end{tabular}
\end{center}
Regras: \textbf{R0} orçamento; \textbf{R1} no máximo 3 investimentos; \textbf{R2} investimento 1 \textbf{ou} 2;
\textbf{R3} \textbf{se} 2 \textbf{então} 1; \textbf{R4} \textbf{se} 2, 3 \textbf{ou} 4 \textbf{então} 1; \textbf{R5}
\textbf{se} 1 \textbf{e} 5 \textbf{então} 3.

\subsection*{Traduzindo regra por regra}

\begin{center}
\begin{tabularx}{\linewidth}{lXl}
\toprule
Regra & Lógica & Restrição\\
\midrule
R0 & orçamento & $4x_1+6x_2+5x_3+3x_4+5x_5\le14$\\
R1 & no máximo 3 de 5 & $x_1+x_2+x_3+x_4+x_5\le3$\\
R2 & $P_1\vee P_2$ (inclusivo) & $x_1+x_2\ge1$\\
R3 & $P_2\Rightarrow P_1$ & $x_2\le x_1$\\
R4 & $P_2\Rightarrow P_1$, $P_3\Rightarrow P_1$, $P_4\Rightarrow P_1$ & (a) $x_2+x_3+x_4\le3x_1$ \ ou \ (b) $x_j\le x_1$, $j=2,3,4$\\
R5 & $P_1\wedge P_5\Rightarrow P_3$ ($k=2$) & $x_1+x_5-x_3\le1$\\
\bottomrule
\end{tabularx}
\end{center}

\paragraph{Por que (a) funciona para R4?} Se $x_1=0$, então $x_2+x_3+x_4\le0$ e nenhum dos três pode ser escolhido.
Se $x_1=1$, então $x_2+x_3+x_4\le3$, que sempre vale. O coeficiente 3 é o ``$M$'' (o maior valor possível da
soma).

\begin{codigo}
def modelo_investimentos(desagregada=True):
    prob = pulp.LpProblem("investimentos", pulp.LpMaximize)
    x = {j: pulp.LpVariable(f"x{j}", cat="Binary") for j in J}
    prob += pulp.lpSum(retorno[j]*x[j] for j in J), "retorno"
    prob += pulp.lpSum(custo[j]*x[j] for j in J) <= ORC, "R0_orcamento"
    prob += pulp.lpSum(x[j] for j in J) <= 3, "R1_max3"
    prob += x[1] + x[2] >= 1, "R2_1_ou_2"
    prob += x[2] <= x[1], "R3_2_implica_1"
    if desagregada:                                     # R4 (b): uma por investimento
        for j in (2, 3, 4):
            prob += x[j] <= x[1], f"R4_{j}_implica_1"
    else:                                               # R4 (a): agregada
        prob += x[2] + x[3] + x[4] <= 3*x[1], "R4_agregada"
    prob += x[1] + x[5] - x[3] <= 1, "R5_1e5_implica_3"
    return prob, x
\end{codigo}
Note que o dicionário \py{x} é indexado de 1 a 5 (\py{J = range(1, 6)}), igual ao enunciado.

\paragraph{Resultado.} Ótimo \textbf{21}, com os investimentos $\{1,2,4\}$ (custo $4+6+3=13$). Sem R5, o ótimo seria
25 com $\{1,4,5\}$; R5 exige o 3 junto com 1 e 5, e com o 3 seriam 4 investimentos, o que viola R1.

\begin{atencao}[diagnóstico pelo resultado]
\begin{itemize}
  \item Deu \textbf{26} com $\{2,4,5\}$? Faltam R3/R4: o investimento 2 está sem o 1.
  \item Deu \textbf{25} com $\{1,4,5\}$? Falta (ou está errada) a R5.
  \item Escreveu R5 como $x_1+x_5\le x_3$? Errado: com $x_1=1,\ x_5=0$, isso exige $x_3=1$. A versão certa é
  $x_1+x_5-x_3\le1$.
\end{itemize}
\end{atencao}

\section{Experimento: agregada $\times$ desagregada}

As duas versões de R4 têm \textbf{exatamente as mesmas soluções inteiras}. Qual é ``melhor''?

\begin{center}
\begin{tabular}{lcc}
\toprule
& valor da relaxação & solução da relaxação\\
\midrule
ótimo inteiro & 21 & $\{1,2,4\}$\\
agregada (a) & 23,71 & $x=(0{,}57;\ 0{,}43;\ 0{,}29;\ 1;\ 0{,}71)$\\
desagregada (b) & 22,5 & $x=(0{,}5;\ 0{,}5;\ 0{,}5;\ 0{,}5;\ 1)$\\
\bottomrule
\end{tabular}
\end{center}

\paragraph{Explicação.} Somando as três restrições de (b), $x_2\le x_1$, $x_3\le x_1$ e $x_4\le x_1$, obtemos
$x_2+x_3+x_4\le3x_1$, que é (a). Logo \emph{todo} ponto que satisfaz (b) satisfaz (a): o poliedro de (b) está
\textbf{dentro} do de (a), e (b) é uma formulação \textbf{mais forte} (Seção~\ref{sec:formulacoes}). A recíproca é
falsa: o ponto $x_1=\tfrac13$, $x_2=1$, $x_3=x_4=0$ satisfaz (a) ($1\le1$) mas viola (b) ($1\not\le\tfrac13$).
Com $x_1$ pequeno, (a) ainda deixa um dos outros inteiro; (b) amarra cada um individualmente.

\begin{resumo}[Lições do Bloco 3]
\begin{itemize}
  \item Menor limite superior = relaxação mais forte = menos nós no \emph{branch-and-bound}.
  \item A formulação desagregada tem \emph{mais} restrições e mesmo assim é \emph{melhor}: ``menos restrições''
  não significa ``modelo melhor''.
  \item Com (b), a R3 ($x_2\le x_1$) fica \textbf{redundante}, porque já aparece em R4. Com (a), ela ainda fortalece a
  relaxação.
\end{itemize}
\end{resumo}
